\documentclass[12pt,a4paper,UTF8]{ctexart}
\usepackage[a4paper,top=22mm,bottom=22mm,left=22mm,right=22mm,headheight=15pt]{geometry}
\newcommand{\Needspace}[1]{\par\begingroup\dimen0=#1\relax\dimen2=\pagegoal\advance\dimen2 by-\pagetotal\ifdim\dimen0>\dimen2\newpage\fi\endgroup}
\newcommand{\needspace}[1]{\par\begingroup\skip0=\lastskip\vskip-\skip0\dimen0=#1\relax\advance\dimen0 by\skip0\vskip\dimen0\penalty-100\vskip-\dimen0\vskip\skip0\penalty9999\endgroup}
\usepackage{amsmath,amssymb,mathtools,bm}
\usepackage{graphicx,xcolor,booktabs,longtable,array,enumitem,fancyhdr,caption,float,microtype,setspace}
\usepackage[hidelinks,unicode]{hyperref}
\usepackage[most]{tcolorbox}
\definecolor{ContestBlue}{HTML}{17365D}
\hypersetup{pdftitle={第070套},pdfauthor={}}
\setstretch{1.24}
\setlength{\parindent}{2em}
\setlength{\parskip}{0.35em}
\setlength{\emergencystretch}{2em}
\setlist{itemsep=0.35em,topsep=0.4em}
\captionsetup{font=small,labelfont=bf,justification=centering}
\setcounter{secnumdepth}{0}
\pagestyle{fancy}
\fancyhf{}
\fancyhead[L]{\small 第070套}
\fancyfoot[C]{\thepage}
\renewcommand{\headrulewidth}{0.4pt}
\ctexset{
  section={format=\Large\bfseries\color{ContestBlue},beforeskip=1.2em,afterskip=0.7em},
  subsection={format=\large\bfseries\color{ContestBlue},beforeskip=1.0em,afterskip=0.5em},
  subsubsection={format=\normalsize\bfseries,beforeskip=0.8em,afterskip=0.35em}
}
\usepackage{calc}
\newcounter{none}
\let\originalsubsection\subsection
\renewcommand{\subsection}{\needspace{6\baselineskip}\originalsubsection}
\makeatletter
\newcommand{\tightlist}{\setlength{\itemsep}{0pt}\setlength{\parskip}{0pt}}
\newcommand{\maxwidth}{\ifdim\Gin@nat@width>\linewidth\linewidth\else\Gin@nat@width\fi}
\newcommand{\maxheight}{\ifdim\Gin@nat@height>0.72\textheight 0.72\textheight\else\Gin@nat@height\fi}
\newcommand{\pandocbounded}[1]{\begingroup\setkeys{Gin}{width=\maxwidth,height=\maxheight,keepaspectratio}#1\endgroup}
\makeatother
\begin{document}
\begin{titlepage}
  \centering
  \vspace*{58mm}
  {\Huge\bfseries\color{ContestBlue} 第070套\par}
  \vfill
  \vspace*{28mm}
\end{titlepage}
\section{第 1 题}\label{ux7b2c-1-ux9898}

水平地面上有一足够长、质量为 \(m\) 的厚平板。平板侧边装有长度为 \(2l\)
的光滑发射台，发射方向与竖直方向夹角为 \(\theta\)。质量为 \(m\)
的小物块由理想弹簧推动；弹簧劲度系数

\[
k=\frac{2mg}{l},
\]

原长为 \(2l\)，压缩至长度 \(l\)
后用细线固定。系统初始静止，剪断细线后小物块沿发射台离开，并最终落回平板上，发射距离记为
\(s\)。

（1）平板固定时，求最大发射距离 \(s_1\) 与相应角度 \(\theta_1\)。

（2）平板可在光滑地面上运动时，最优角为 \(\theta_2\)。判断
\(\theta_2-\theta_1\) 的正负。

\begin{figure}[H]
\centering
\includegraphics[width=0.8\linewidth,height=\textheight,keepaspectratio,alt={图1-1　弹簧发射台与厚平板}]{../figures/question-01.pdf}
\caption*{图1-1　弹簧发射台与厚平板}
\end{figure}

\newpage

\section{第 2 题}\label{ux7b2c-2-ux9898}

空间充满相对介电常数为 \(\varepsilon_{r2}\) 的均匀介质，远处匀强电场为

\[
\boldsymbol E_\infty=(E_0,0,E_0).
\]

在介质中挖出半径为 \(r\)、轴线沿 \(z\)
轴的无限长圆柱形孔，并在孔内填入相对介电常数为 \(\varepsilon_{r1}\)
的介质。求电场分布。

\begin{figure}[H]
\centering
\includegraphics[width=0.487\linewidth,height=\textheight,keepaspectratio,alt={图2-1　介质圆柱嵌入体与斜向外电场}]{../figures/question-02.pdf}
\caption*{图2-1　介质圆柱嵌入体与斜向外电场}
\end{figure}

\newpage

\section{第 3 题}\label{ux7b2c-3-ux9898}

椭球壳内表面满足

\[
\frac{x^2}{a^2}+\frac{y^2}{b^2}+\frac{z^2}{c^2}=1.
\]

底部有半径为 \(r\)、质量为 \(m\)
的均匀实心小球，摩擦足够大，始终纯滚动。

（1）椭球壳固定。\(t=0\) 时球心位于最低点附近 \((0,y_0,z_0)\)，速度为
\((v_0,0,0)\)，并满足

\[
y_0\ll a,b,c,r,
\qquad
\frac{v_0^2}{g}\ll a,b,c,r.
\]

在小振幅近似下求 \(x(t)\)、\(y(t)\)。

（2）椭球壳以角速度 \(\omega_0\) 绕 \(z\)
轴匀速转动，求最低点附近小幅运动的简正频率。

\begin{figure}[H]
\centering
\includegraphics[width=0.636\linewidth,height=\textheight,keepaspectratio,alt={图3-1　椭球壳内滚动小球与三维坐标}]{../figures/question-03.pdf}
\caption*{图3-1　椭球壳内滚动小球与三维坐标}
\end{figure}

\newpage

\section{第 4 题}\label{ux7b2c-4-ux9898}

如图所示，自制手推炮的炮膛截面积为 \(S\)，工作过程分为以下四个状态。

\begin{figure}[H]
\centering
\includegraphics[width=0.45\linewidth,height=\textheight,keepaspectratio,alt={图4-1　手推炮的四个工作状态与气柱几何量}]{../figures/question-04.pdf}
\caption*{图4-1　手推炮的四个工作状态与气柱几何量}
\end{figure}

\begin{enumerate}
\def\labelenumi{\arabic{enumi}.}
\tightlist
\item
  炮膛中已有一颗炮弹 \(A\)，再从炮尾塞入一颗与其相同的炮弹
  \(B\)。两弹质量均为
  \(m\)，体积均可忽略，且与炮膛之间保持良好气密性。随后以速度 \(u\)
  匀速推动手柄，压缩两弹之间的气体。
\item
  当炮弹 \(B\) 被推到距炮口 \(l_1+l_2\) 处时，炮弹 \(A\) 已运动到距炮口
  \(l_1\) 处，并获得速度 \(u\)。此时两弹间气体压强为
\end{enumerate}

\[
P_0=\frac{ma_0}{S}.
\]

此后 \(B\) 继续由手柄匀速推进，\(A\)
在气体压力作用下加速。作如下理想化：当炮弹相对炮膛的速度不小于 \(u\)
时，炮弹与炮膛之间的摩擦可忽略；气体经历绝热膨胀，且其压强始终远大于大气压；绝热指数为

\[
   \gamma=\frac{C_P}{C_V}=2.
   \]

\begin{enumerate}
\def\labelenumi{\arabic{enumi}.}
\setcounter{enumi}{2}
\tightlist
\item
  从状态 2 到状态 3，炮弹 \(B\) 被手柄匀速推进距离 \(h\)，炮弹 \(A\)
  被加速至速度 \(v\) 并射出。
\item
  发射完成后，炮弹 \(B\) 继续向炮口推进并成为下一次发射中的炮弹 \(A\)。
\end{enumerate}

（1）忽略状态 2 到状态 3 中炮弹 \(B\) 的推进距离 \(h\)，求炮弹 \(A\)
的发射速度 \(v_1\)，结果用 \(u,l_1,l_2,a_0\) 表示。

（2）计入推进距离 \(h\)，求炮弹 \(A\) 的发射速度 \(v_2\)，结果用
\(h,u,l_1,l_2,a_0\) 表示。

（3）根据炮弹的动力学过程确定 \(h\) 与 \(u,l_1,l_2,a_0\) 的关系。若
\(h=l\)、\(l_1=2l\)、\(l_2=3l\)，求 \(u\)，结果用 \(l,a_0\) 表示。

\newpage

\section{第 5 题}\label{ux7b2c-5-ux9898}

已知四边形 \(ABCD\) 是一个正方形物体通过理想薄凸透镜所成的像。
请根据该四边形，用尺规作出原正方形和透镜的位置，并写出作图的关键步骤。

\begin{figure}[H]
\centering
\includegraphics[width=0.767\linewidth,height=\textheight,keepaspectratio,alt={图5-1　由同一薄透镜公式生成的正方形像四边形}]{../figures/question-05.pdf}
\caption*{图5-1　由同一薄透镜公式生成的正方形像四边形}
\end{figure}

\newpage

\section{第 6 题}\label{ux7b2c-6-ux9898}

一根长为 \(l\)、质量为 \(m\) 的均匀刚性细杆在桌边一个很细的固定凸起上方
\(h\)
处保持水平，并由静止释放。细杆与凸起发生完全非弹性碰撞，碰撞点为细杆右侧三等分点。

\begin{enumerate}
\def\labelenumi{\arabic{enumi}.}
\tightlist
\item
  求碰撞后瞬间细杆质心速度与绕质心转动的角速度。
\item
  杆与凸起间的摩擦因数为
  \(\mu=2\)。若碰撞后瞬间接触点不相对滑动，求允许的最大释放高度
  \(h_{\max}\)。
\item
  取
  \(h=h_{\max}/2\)，求碰撞后杆与凸起刚要发生相对滑动时，杆与水平方向的夹角
  \(\theta_0\)。
\item
  写出发生相对滑动后细杆的运动微分方程。
\end{enumerate}

\begin{figure}[H]
\centering
\includegraphics[width=0.621\linewidth,height=\textheight,keepaspectratio,alt={图6-1　水平细杆自由下落并与桌边凸起碰撞}]{../figures/question-06.pdf}
\caption*{图6-1　水平细杆自由下落并与桌边凸起碰撞}
\end{figure}

\newpage

\section{第 7 题}\label{ux7b2c-7-ux9898}

顶角为 \(45^\circ\) 的金属导轨 \(MON\)
固定在水平面内，并处于竖直方向、磁感应强度为 \(B\)
的匀强磁场中。一根始终垂直于导轨 \(ON\)
的导体棒与两导轨接触，沿水平方向向右运动。导体棒质量为
\(m\)，导轨与导体棒单位长度的电阻均为 \(r_0\)，接触电阻忽略。以顶点
\(O\) 为坐标原点，棒的位置坐标为 \(x\)。

\begin{enumerate}
\def\labelenumi{\arabic{enumi}.}
\tightlist
\item
  棒以恒定速度 \(v_0\) 运动时，求外力 \(F\)。
\item
  求从 \(0\) 到 \(t\) 时间内产生的焦耳热 \(Q(t)\)。
\item
  在 \(t=t_0\) 时撤去外力，求棒最终位置 \(x_f\)。
\item
  若棒从 \(O\) 点以初速度 \(v_1\) 出发，此后不受外力，求最终位置及
  \(x(t)\)。
\end{enumerate}

\begin{figure}[H]
\centering
\includegraphics[width=0.447\linewidth,height=\textheight,keepaspectratio,alt={图7-1　匀强磁场中的直角三角形变长度导轨回路}]{../figures/question-07.pdf}
\caption*{图7-1　匀强磁场中的直角三角形变长度导轨回路}
\end{figure}

\newpage

\section{第 8 题}\label{ux7b2c-8-ux9898}

存在连续衰变链

\[
{}^{228}\mathrm{Th}\longrightarrow{}^{224}\mathrm{Ra}
\longrightarrow{}^{220}\mathrm{Rn}\longrightarrow\cdots,
\]

前两次衰变的半衰期分别为 \(\tau_1=1.91\,\mathrm{a}\) 和
\(\tau_2=3.66\,\mathrm d\)。某矿物中三种核素已达到放射性平衡。从中提取全部
\({}^{228}\mathrm{Th}\) 与 \({}^{224}\mathrm{Ra}\)，组成质量
\(M=1\,\mathrm g\) 的混合物。求：

\begin{enumerate}
\def\labelenumi{\arabic{enumi}.}
\tightlist
\item
  \(t=0\) 时两种核素的核数 \(N_{10},N_{20}\)；
\item
  \(t>0\) 时 \({}^{224}\mathrm{Ra}\) 的核数 \(N_2(t)\)；
\item
  当 \({}^{228}\mathrm{Th}\) 减为初值一半时，\({}^{224}\mathrm{Ra}\)
  的核数。
\end{enumerate}

\begin{figure}[H]
\centering
\includegraphics[width=0.88\linewidth,height=\textheight,keepaspectratio,alt={图8-1　钍、镭、氡的两级连续放射性衰变链}]{../figures/question-08.pdf}
\caption*{图8-1　钍、镭、氡的两级连续放射性衰变链}
\end{figure}
\end{document}
